\subsection{Which conclusions change after projection?}\label{sec:decision-application}
The primary question is which measurements can be certified as leaves, excluded from hub status, or placed within the top three. These are statements about the marginal association tree; leaf status alone does not justify removing a measurement from a predictive model. Table~\ref{tab:decision-abalone} separates the gain from replacing one matching by E64 from the gain due to subsequent optimization.


\begin{table}[H]\centering\small
\setlength{\tabcolsep}{4pt}
\caption{Abalone pooled: concrete conclusions under the IID working model and a separately chosen calibration. All rows use local edge comparisons. Fixed and excluded denote counts among eight measurements. The last two columns concern whole weight. P is one-matching exact McNemar--Holm. The Full LP rows use path projection. Exact enumeration (main text, Section~\ref{sec:sharp}) further improves E64 Height degree from $[1,3]$ to $[1,2]$ and raises hub exclusions from four to five.}\label{tab:decision-abalone}
\begin{lrbox}{\resultsbox}\begin{tabular}{llrrrr}
\toprule
Test & Projection & Fixed & Excluded & Degree & Rank\\
\midrule
P & Full LP & 1 & 1 & [1,6] & [1,8]\\
E64 & Minor & 1 & 2 & [2,5] & [1,6]\\
E64 & Degree LP & 2 & 4 & [3,5] & [1,3]\\
E64 & Full LP & 2 & 4 & [3,5] & [1,3]\\
\bottomrule\end{tabular}\end{lrbox}\ifdim\wd\resultsbox>\linewidth\resizebox{\linewidth}{!}{\usebox{\resultsbox}}\else\usebox{\resultsbox}\fi\end{table}


With E64 fixed, the minor stage excludes length and shucked weight from hub status. Degree-only path projection additionally excludes viscera weight and ring count, certifies ring count as a leaf, and raises the lower degree of whole weight from two to three. Together with degree conservation, these statements place whole weight within the top three. It is certified to have strictly greater degree than length, shucked weight, viscera weight and ring count in every population-optimal tree. The extra 56 ordered-pair LP objectives change none of the reported degree or rank endpoints. The top-three statement has a simple structural explanation. The eight degrees sum to 14. If whole weight and three competitors each had degree at least three, these four vertices plus the four remaining positive degrees would sum to at least 16. At most two competitors can therefore match or exceed whole weight. The inferential conclusion is four excluded hubs and a top-three restriction, with the identity of a unique population hub still unresolved.

Table~\ref{tab:decision-populations} reports the entire E64 local-family measurement grid, including cases with no new decision. At $n=1307$ and $1528$, degree LP certifies ring count as a leaf in the female and male Abalone strata, respectively; neither minor projection resolves those degrees. In the infant stratum, shucked weight is already a certified leaf before LP. The corresponding guarantees require an IID model separately within each analysed population. The comparisons across populations, methods and families are descriptive and have no joint 95\% coverage statement.


\begin{table}[H]\centering\small
\setlength{\tabcolsep}{4pt}
\caption{All 12 measurement populations: E64 local-family decision counts, minor $\to$ degree-only LP. Full ordered-pair projection gives the same rank endpoints as degree-only LP in every row. Dry Bean remains a computational stress study because the sampling design and repeated rows limit a population interpretation.}\label{tab:decision-populations}
\begin{lrbox}{\resultsbox}\begin{tabular}{lrrrrr}
\toprule
Population & $n$ & $p$ & Fixed & Excluded & Top three\\
\midrule
Abalone pooled & 4177 & 8 & 1 $\to$ 2 & 2 $\to$ 4 & 0 $\to$ 1\\
Abalone F & 1307 & 8 & 0 $\to$ 1 & 0 $\to$ 1 & 0 $\to$ 0\\
Abalone I & 1342 & 8 & 1 $\to$ 1 & 1 $\to$ 1 & 0 $\to$ 0\\
Abalone M & 1528 & 8 & 0 $\to$ 1 & 0 $\to$ 1 & 0 $\to$ 0\\
Bean Barbunya & 1322 & 16 & 0 $\to$ 0 & 0 $\to$ 0 & 0 $\to$ 0\\
Bean Bombay & 522 & 16 & 0 $\to$ 0 & 0 $\to$ 0 & 0 $\to$ 0\\
Bean Cali & 1630 & 16 & 0 $\to$ 0 & 0 $\to$ 0 & 0 $\to$ 0\\
Bean Dermason & 3546 & 16 & 0 $\to$ 0 & 0 $\to$ 0 & 0 $\to$ 0\\
Bean Horoz & 1928 & 16 & 0 $\to$ 0 & 0 $\to$ 0 & 0 $\to$ 0\\
Bean Seker & 2027 & 16 & 0 $\to$ 0 & 0 $\to$ 0 & 0 $\to$ 0\\
Bean Sira & 2636 & 16 & 0 $\to$ 0 & 0 $\to$ 0 & 0 $\to$ 0\\
Bean Pooled & 13611 & 16 & 0 $\to$ 0 & 0 $\to$ 0 & 0 $\to$ 0\\
\bottomrule\end{tabular}\end{lrbox}\ifdim\wd\resultsbox>\linewidth\resizebox{\linewidth}{!}{\usebox{\resultsbox}}\else\usebox{\resultsbox}\fi\end{table}


Across all 192 original method--family--population analyses, degree-only LP adds at least one of the stated decisions in 32 analyses. Full pair projection gives exactly the same rank endpoints in 188 of 192 analyses. Its remaining benefit is six additional hub exclusions across four analyses: female Abalone with multiplier rectangles and male Abalone with P, each under both comparison families. These four analyses concern two populations and are not independent confirmations. No stage produces a new unique-hub declaration on the measurement grid. Dry Bean's E64 degree widths narrow, but no leaf, hub exclusion or top-three restriction is certified, even in the pooled sample of 13,611 records. This distinguishes a narrower numerical bound from a new substantive conclusion.
